% ===
\section{Introdução}
% ===

% ------------------
%%% PANORAMA
% ------------------

%% Contextualização geral
A relação entre a atividade econômica e o meio ambiente natural tem sido há muito marcada por uma tensão persistente entre expansão produtiva e preservação ecológica. Muito embora o meio ambiente é essencial ao bem-estar humano e à produção econômica, o esgotamento dos recursos naturais e a degradação dos sistemas ecológicos podem comprometer a própria capacidade produtiva que a atividade econômica procura expandir. Neste sentido, produção, investimento e progresso tecnológico geram emprego, renda e melhorias nos padrões de vida, de modo que a política ambiental não pode ser avaliada independentemente de seus efeitos sobre os incentivos privados e a viabilidade produtiva. Portanto, o problema econômico relevante não é se a produção ou a preservação deve prevalecer, mas sob quais condições o uso da natureza pode permanecer compatível com as funções ecológicas das quais dependem a produção presente e futura.

%% Teoria atual
A teoria econômica não ignora a natureza. Terra, recursos naturais, custos ambientais, externalidades e recursos de uso comum têm sido amplamente examinados em diferentes vertentes da literatura, como em \citeonline{pigou2010liberty}, \citeonline{coase1960social}, \citeonline{hardin1968tragedy}, \citeonline{ostrom1990commons} e \citeonline{dasgupta2001human}, conforme veremos detalhadamente na Seção 3.

Nas representações canônicas da produção, como a função Cobb–Douglas, o produto é geralmente modelado como uma função do capital e do trabalho, por vezes com a terra ou os recursos naturais figurando como insumos, mas os sistemas ecológicos não são, em geral, representados como um fator ou cofator que afeta o produto, a produtividade futura, os incentivos estratégicos e o bem-estar intertemporal. Em vez disso, o meio ambiente frequentemente aparece como um estoque, restrição, externalidade ou objeto de regulação corretiva, e essa representação, por sua vez, produz consequências econômicas. Ato contínuo, essa escolha de modelagem não é neutra: ao representar o meio ambiente principalmente como estoque, restrição ou externalidade, os modelos convencionais limitam a medida em que a funcionalidade ecológica pode ingressar nos problemas de decisão privada como um estado que codetermina a produtividade.

%% Implicações

% Teóricas
Quando a funcionalidade ecológica permanece fora do problema de decisão privada, a preservação tende a aparecer como uma obrigação legal, um custo de conformidade ou um custo de oportunidade imposto à produção. Seus benefícios são frequentemente postergados, incertos, espacialmente difusos ou apropriados coletivamente, enquanto os retornos da conversão de recursos podem ser imediatos, mensuráveis e apropriados privadamente. Portanto, mesmo um agente que reconheça a importância de longo prazo da integridade ecológica pode racionalmente escolher a expansão de curto prazo quando o retorno privado esperado da degradação supera o retorno privado esperado da preservação. Essa divergência é intensificada por razões conjunturais diversas, como restrições de crédito, pressões competitivas, informação incompleta, fiscalização deficiente e incerteza quanto às políticas ou às condições de mercado futuras.

% Práticas
A regulação governamental de comando e controle, inclusive com penalidades várias como multas, embargos judiciais e ditames burocráticos, mercados de carbono, certificações ambientais, pagamentos por serviços ecossistêmicos e incentivos fiscais traduzem partes dessas percepções em políticas públicas. Não obstante ser possível asseverar que essas abordagens produziram melhorias importantes em dimensões específicas, não é descabido sopesar que elas tratam de componentes diferentes do problema e não estabelecem automaticamente uma relação unificada entre condições ecológicas, produção, retornos privados e governança territorialmente fundamentada.

%% Lacunas (teóricas e práticas)
Dito isso, a lacuna remanescente não consiste nem na ausência da natureza no pensamento econômico nem na ausência de instrumentos legais, financeiros, institucionais ou normativos destinados a protegê-la. O que permanece insuficientemente especificado é, antes, um conjunto integrado de condições sob as quais a funcionalidade ecológica seja conjuntamente determinada com a produção, seus efeitos produtivos e relativos aos serviços ecossistêmicos ingressem nos retornos privados e a governança corresponda à escala espacial dos processos ecológicos relevantes. Assim, faz-se mister registrar, a lacuna abordada neste artigo é de integração analítica, e não de ausência de considerações ambientais na economia.

% ------------------
%%% EPISTEME
% ------------------

%% PROBLEMA
Diante desse contexto, o artigo busca resolver o seguinte \textbf{problema}: \emph{sob quais condições a funcionalidade ecológica pode ser incorporada como um estado endógeno da produção, e quais arranjos institucionais podem tornar sua manutenção compatível, em termos de incentivos, com a viabilidade econômica da atividade privada?}

Essa formulação substitui a oposição simplificada entre uso e preservação por um problema de desenho produtivo e institucional. O foco deixa de ser apenas a existência de impactos ambientais decorrentes da atividade econômica e passa a recair sobre uma questão mais específica: se essa atividade degrada sistematicamente as condições ecológicas das quais dependem a produção futura e o bem-estar social.

%% OBJETIVOS
O \textbf{objetivo geral} do artigo é desenvolver um arcabouço teórico-analítico que represente o meio ambiente como um fator endógeno da produção, cuja condição é simultaneamente afetada pelas decisões produtivas e capaz de influenciar seus resultados. A partir desse arcabouço, o artigo busca identificar as condições institucionais e econômicas sob as quais a manutenção desse fator se torna compatível com os incentivos e a viabilidade econômica dos agentes privados, bem como definir e caracterizar o comensalismo econômico como uma configuração ecológico-econômica distinta.

Mais especificamente, os \textbf{objetivos específicos} são: (i) distinguir o uso de recursos naturais como insumos convencionais da representação das condições ambientais relevantes como um fator endógeno que afeta a produção e é, simultaneamente, afetado pelas decisões de uso e preservação; (ii) examinar como a heterogeneidade espacial e a substituibilidade imperfeita dos serviços ecossistêmicos condicionam a escala territorial da governança ambiental e dos mecanismos de compensação; (iii) derivar as condições de incentivos e retorno privado sob as quais a manutenção do fator ambiental se torna uma escolha individualmente racional para os agentes produtivos; e (iv) definir formalmente o comensalismo econômico e comparar suas implicações produtivas e ambientais com aquelas associadas aos regimes predatório, extrativista e mutualista de uso dos recursos.

%% HIPÓTESES

% Geral
A análise é orientada pela seguinte \textbf{hipótese geral}: quando as condições ambientais relevantes são representadas como um fator endógeno que influencia os resultados produtivos e é, simultaneamente, afetado pelas decisões de uso e preservação, arranjos institucionais territorialmente ajustados podem incorporar os benefícios e custos ambientais ao cálculo econômico dos agentes privados. Esse alinhamento de incentivos pode tornar a viabilidade econômica da atividade produtiva compatível com a manutenção das condições ambientais das quais ela depende. O artigo denomina essa possibilidade de configuração comensal e apresenta sua definição formal nas seções seguintes. % definir depois a distinção de comensalismo de sustentabilidade. PAREI AQUI

% Específicas
Essa hipótese geral é decomposta em três \textbf{hipóteses específicas}.

(H1) Endogeneidade produtivo-ecológica. Quando a funcionalidade ecológica afeta o produto e, simultaneamente, responde à produção, à preservação e às regras institucionais, tratá-la como um insumo fixo ou exclusivamente como um efeito externo omite mecanismos de retroalimentação que alteram os retornos marginais e podem modificar a ordenação relativa das estratégias alternativas de uso dos recursos. % Modifica as condições marginais de produção.

(H2) Governança local. Quando as funções ecossistêmicas são espacialmente heterogêneas e substitutas imperfeitas, mecanismos de governança e compensação ancorados em unidades territoriais ecologicamente relevantes podem reduzir a divergência entre os retornos privados e territoriais da preservação de maneira mais eficaz do que mecanismos espacialmente fungíveis, desde que os efeitos ecológicos relevantes sejam observáveis, as instituições locais possuam capacidade suficiente de implementação e os efeitos transfronteiriços sejam coordenados em níveis institucionais mais abrangentes. % Modifica a escala e a incidência institucional.

(H3) Compatibilidade intertemporal de incentivos. Quando os retornos privados incorporam a contribuição da funcionalidade ecológica para a produção corrente e futura esperada, os serviços ecossistêmicos remunerados e as perdas evitadas decorrentes da degradação, amplia-se o intervalo de custos de oportunidade para o qual a preservação é privadamente ótima, em comparação com um retorno baseado exclusivamente no produto corrente comercializado. A preservação torna-se individualmente racional sempre que os benefícios marginais privados resultantes igualam ou superam seus custos de oportunidade e de implementação. % Modifica os incentivos intertemporais dos agentes.

Essas hipóteses são analíticas, e não diretamente empíricas. H1 é examinada mediante a comparação das condições marginais obtidas sob representações exógenas e endógenas do estado ecológico; H2, mediante a comparação de configurações espaciais e institucionais alternativas; e H3, mediante a derivação do limiar no qual o retorno privado da manutenção da funcionalidade ecológica se torna pelo menos tão elevado quanto o retorno da degradação. A estabilidade dinâmica dessa configuração sob incerteza, aprendizagem e decisões repetidas é deixada para o modelo intertemporal subsequente.

%% Contribuições
O artigo oferece três contribuições principais. Primeiro, distingue a inclusão dos recursos naturais como insumos convencionais da endogeneidade da funcionalidade ecológica como um estado da produção conjuntamente determinado. Segundo, conecta essa representação produtiva aos retornos privados e aos arranjos de governança ancorados em unidades territoriais ecologicamente relevantes. Terceiro, define o comensalismo econômico por meio do comportamento conjunto da viabilidade privada e da funcionalidade ecológica e deriva as condições sob as quais essa configuração pode emergir em relação aos regimes predatório, extrativista e mutualista. Em conjunto, essas contribuições fornecem uma base conceitual e analítica para tratar a preservação não apenas como uma restrição externa, mas como um componente potencialmente compatível, em termos de incentivos, com as estratégias produtivas.

%% Metodologia
Metodologicamente, o artigo adota uma abordagem teórico-analítica que combina teoria da produção, economia ambiental e análise institucional. Primeiro, reconstrói como as externalidades, os direitos de propriedade, os recursos de uso comum e o capital natural esclarecem diferentes partes do problema produção-preservação. Em seguida, desenvolve um modelo preliminar no qual o estado ambiental ingressa na estrutura de produção e no retorno dos agentes produtivos, enquanto os incentivos e a governança afetam o retorno relativo da preservação e da exploração. O modelo é utilizado para comparar estratégias ecológico-econômicas alternativas e derivar as condições sob as quais uma configuração comensal pode emergir como um resultado individualmente racional e ecologicamente não degradante. A análise permanece deliberadamente fundacional: estabelece a base analítica conceitual e estática sobre a qual a dinâmica intertemporal, a incerteza, a aprendizagem e os testes empíricos poderão ser posteriormente desenvolvidos.

%% Divisão do artigo
O restante do artigo está organizado da seguinte forma. A Seção 2 contextualiza o problema produção-preservação e as limitações da governança ambiental atual. A Seção 3 revisa a literatura relevante, partindo das externalidades e dos direitos de propriedade em direção à governança dos recursos de uso comum e ao capital natural, e identifica os elementos que permanecem analiticamente desconectados. A Seção 4 desenvolve o modelo conceitual e analítico do comensalismo econômico, define os regimes ecológico-econômicos alternativos e estabelece as condições de equilíbrio relevantes. A Seção 5 apresenta os resultados teóricos e as implicações comparativas do modelo. A Seção 6 discute sua relação com a literatura existente, seu significado institucional e as limitações do arcabouço proposto. A Seção 7 conclui e indica caminhos para a formalização dinâmica e a avaliação empírica.

% fim %
